\documentclass[11pt,a4paper]{article}
\usepackage[T1]{fontenc}
\usepackage{isabelle,isabellesym}
\usepackage[a4paper,margin=2.4cm]{geometry}
\usepackage{amsmath}
\usepackage{amssymb}
\usepackage{booktabs}
\usepackage{microtype}
\usepackage{url}
\usepackage[hidelinks,bookmarksopen=true]{hyperref}

\urlstyle{rm}
\isabellestyle{it}
\setlength{\emergencystretch}{3em}

\begin{document}

\title{Foundations and Limits of Essentially Ordered Causal Regress Arguments}
\author{Arthur F. Ramos \and David Barros Hulak \and Ruy J. G. B. de Queiroz}
\maketitle

\begin{abstract}
This entry formalizes a carrier-restricted theory of causal and grounding
regress arguments in Isabelle/HOL.  It distinguishes acyclicity,
well-foundedness, explicit chain termination, local foundation, global
foundation, and full grounding, and proves the implications that hold between
them together with finite and infinite separation models.  An executable finite
checker returns either an exact independent-root certificate or a concrete
missing-foundation witness, with kernel-checked correctness theorems.  The
entry also formalizes a bounded cycle analyzer with explicit finite return-path
certificates,
a certified finite profile linking five declarative properties to two executable
tests, and premise-deletion countermodels
for the Romero--Pérez and Cohoe-style packages.  It formalizes two explicit
Second-Way reconstructions: a Cohoe-style reconstruction in which an
essentially ordered-series explication, explicit foundation, and bridge
conditions yield a non-derivative first cause, and a Romero--Pérez-style
  well-founded reconstruction with its exact independent-member conclusion.  The
  results are conditional on stated dependency premises; no theorem identifies a
  formal first cause with God, and the substantive application is deliberately
  focused on the Second Way.
\end{abstract}

\section{Contribution and scope}

The phrase \lq\lq there must be a first cause'' compresses several distinct
claims.  Acyclicity need not be well-foundedness; a selected target may have a
foundation while the whole relation is non-terminating; and global foundation
uses a set of independent supporters, not a singleton first-cause claim.  This
entry makes those distinctions explicit and mechanically checkable.

Its contribution is a compact package of four layers:

\begin{enumerate}
\item a carrier-restricted interface for dependence, ancestry, grounding,
  independence, and first causes;
\item an implication/separation atlas for acyclicity, termination, and
  local/global foundation;
\item executable finite certificates for independent roots, missing
  foundations, and closed cycle paths; and
\item source-oriented Second-Way formal explications with premise-deletion
  countermodels.
\end{enumerate}

All conclusions are conditional on stated relation-theoretic assumptions.  The
entry does not assert that a physical or metaphysical domain is well-founded,
or identify any formal first cause with God.  It deliberately does not attempt
the First or Third Ways: substantive treatments would require richer
act/potency and modal semantics beyond this relation-theoretic scope.

\section{Formal interface}

Fix a type of entities, a carrier \(C\), and an immediate-dependence relation
\(D\).  The relation is read as \lq\lq \(x\) is an immediate essential cause,
ground, or actualizer of \(y\).''  The project uses carrier-restricted
predicates:

\[
  D_C(x,y) \mathrel{\equiv} x\in C \land y\in C \land D\,x\,y.
\]

The predicate \texttt{dependency\_ancestor} is the carrier-restricted
reflexive-transitive closure of immediate dependence, while
\texttt{grounded\_in} records an immediate predecessor.  Independence means
that no such predecessor exists; \texttt{first\_member\_in} and
\texttt{first\_cause\_in} add reachability and the required non-identity
condition.  The interface is neutral about whether \(D\) means efficient
causation, grounding, or actualization.

\section{Foundation principles}

The entry formalizes the following predicates.

\begin{center}
\begin{tabular}{@{}p{.25\linewidth}p{.64\linewidth}@{}}
\toprule
\textbf{Predicate} & \textbf{Formal content} \\
\midrule
\texttt{acyclic\_on} & no nonempty dependency cycle lies in the carrier; \\
\texttt{well\_founded\_on} & every nonempty carrier subset has a minimal member; \\
\texttt{chain\_terminating\_on} & every infinite dependency chain in the carrier is impossible; \\
\texttt{local\_foundation\_on} & every immediately grounded target in a selected set has a reachable independent member; \\
\texttt{global\_foundation\_on} & a designated set of independent members supports every grounded target in a selected set; \\
\texttt{full\_foundation\_on} & every partially grounded target in the selected set has a full ground contained in a common independent foundation set. \\
\bottomrule
\end{tabular}
\end{center}

The local and global predicates are weaker than well-foundedness: they concern
selected targets and do not quantify over every nonempty subset.  In particular,
\texttt{global\_foundation\_on} is witnessed by a set \(F\) of independent
members; it is not a claim that there is one unique common first cause.  Full
grounding remains distinct from partial grounding.

The headline hierarchy is summarized below.  The last two rows are strict
separations, not additional assumptions.

\begin{center}
\small
\begin{tabular}{@{}p{.38\linewidth}p{.52\linewidth}@{}}
\toprule
\textbf{Checked relation} & \textbf{Scope and witness} \\
\midrule
\(\mathrm{WF}\Longleftrightarrow\mathrm{CT}\) & unrestricted equivalence of
  well-foundedness and chain termination; \\
\(\mathrm{WF}\Longrightarrow\mathrm{AC}\) & unrestricted implication to
  acyclicity; \\
\(\mathrm{WF}\Longrightarrow\mathrm{LF}\); for \(R\subseteq C\),
  \(\mathrm{LF}\Longleftrightarrow\mathrm{GF}\) &
  well-foundedness implies local foundation without extra coverage; the
  local/global collapse requires carrier coverage; \\
\(\mathrm{AC}\Longleftrightarrow\mathrm{WF}\Longleftrightarrow\mathrm{CT}
  \Longrightarrow\mathrm{LF}\); for \(R\subseteq C\),
  \(\mathrm{LF}\Longleftrightarrow\mathrm{GF}\) &
  finite carrier; the rooted-cycle model refutes the final converse; \\
\(\mathrm{AC}\not\Longrightarrow\mathrm{WF},\mathrm{LF}\) & integer
  predecessor model; \\
\(\mathrm{LF},\mathrm{GF}\not\Longrightarrow\mathrm{AC}\) & finite
  rooted-cycle model. \\
\bottomrule
\end{tabular}
\end{center}

\section{Implication and separation atlas}

The theorem \texttt{foundation\_implication\_atlas} collects the positive
implications and names the bridge assumptions needed to convert full grounding
into a local first-cause conclusion: target presence, target-set nonemptiness,
and ancestry coverage.  None is hidden in a simp rule.

\subsection{Infinite separation models}

The integer predecessor model is acyclic and gives every integer a predecessor,
but has no independent member.  The pedestal model has a foundation set while
remaining non-well-founded.  The finite rooted-cycle model has foundation for
selected targets while retaining a separate cycle.  The paired model separates
partial from full grounding.  These are executable Isabelle/HOL models used by
the Second-Way comparison layer, not informal diagrams.

\section{Executable finite checker}

\texttt{Finite\_Foundation\_Checker.thy} represents a finite model by a
\texttt{nat list} of vertices and ordered edge pairs, then derives the carrier,
relation, bounded reachability, grounding, and independence predicates.  The
path recursion is proved sound and complete under the finite-list invariant.

The executable function \texttt{finite\_foundation\_result} returns one of two
constructors:

\begin{center}
\begin{tabular}{@{}p{.42\linewidth}p{.49\linewidth}@{}}
\toprule
\textbf{Result} & \textbf{Meaning} \\
\midrule
\texttt{Foundation\_Certificate xs} &
\texttt{xs} is exactly the list of independent vertices in the checked
carrier, and its set is proved to be a \texttt{foundation\_set\_for}; \\
\texttt{Missing\_Foundation\_Certificate x} &
\texttt{x} is a concrete carrier vertex that is grounded but has no reachable
independent member. \\
\bottomrule
\end{tabular}
\end{center}

The correctness theorems connect the result with
\texttt{finite\_foundation\_check}, identify the positive payload with the
exact independent-root set, and certify the negative witness.  Regression
models cover a finite chain and directed cycle; the checker is therefore a
small certified assay, not an appeal to an external graph library.

\subsection{Cycle certificates and the finite profile}

\texttt{Finite\_Foundation\_Analysis.thy} adds a second executable result
type.  \texttt{No\_Cycle\_Certificate} is returned exactly when the
carrier-restricted relation is acyclic.  Otherwise
\texttt{Cycle\_Certificate x xs} returns a carrier vertex \(x\) together with a
nonempty finite path list \(xs\) for which the formal transitive closure
contains a return path
\[
  D_C^{+}(x,x).
\]
The certificate theorem checks the closed \texttt{rtrancl\_path} represented by
\(xs\), rather than returning only an endpoint.  The internal search is bounded
by the listed vertices; the emitted path is at most one longer because it
prepends the first edge's endpoint.  This is proved by
\texttt{finite\_cycle\_result\_cycle\_length}.

The exported \texttt{finite\_foundation\_profile} packages five Boolean
answers, computed by two executable tests:
\[
(\text{acyclic},\text{well-founded},\text{chain-terminating},
  \text{local-foundation},\text{global-foundation}).
\]
The theorem \texttt{finite\_foundation\_profile\_correct} proves equality with
the declarative tuple for every finite input model.  Thus the profile exposes
the finite collapse at the executable API boundary while making clear that
acyclicity, well-foundedness, and chain termination share one cycle test, and
local and global foundation share one foundation test.

\section{Two explicit Second-Way reconstructions}

The source comparison is guided by the literature on essentially ordered
causal series.  Cohoe's analysis emphasizes that Aquinas's argument concerns a
series whose later members derive their causal power from earlier members
\cite{Cohoe2013}.  Romero and Pérez present a first-order reconstruction whose
formal source-level conclusion is the existence of an uncaused entity, while
warning against silently adding uniqueness or a universe-cause conclusion
\cite{RomeroPerez2012}.

\subsection{Source-to-formal premise map}

The following table separates source-level commitments from the Isabelle
encoding.  In particular, the bridge conditions are recorded as formal
translation assumptions, not presented as quotations from the philosophical
literature.

\begin{center}
\small
\begin{tabular}{@{}p{.22\linewidth}p{.39\linewidth}p{.29\linewidth}@{}}
\toprule
\textbf{Source focus} & \textbf{Formal counterpart} & \textbf{Role in the result}\\
\midrule
Cohoe: an essentially ordered causal series
  \cite{Cohoe2013} &
\texttt{essentially\_ordered\_series}, with a derivative predicate and a
carrier-restricted immediate relation &
An explicit explication; it is not identified with well-foundedness.\\
Cohoe: foundation of the selected series &
\texttt{full\_foundation\_on} plus endpoint, nonemptiness, and ancestry
bridges &
Supplies the binary foundation conclusion only under stated translation
assumptions.\\
Romero--Pérez: caused items and predecessor coverage
  \cite{RomeroPerez2012} &
Existence, carrier, and predecessor premises in
\texttt{romero\_perez\_second\_way\_premises} &
Feeds the per-target first-cause theorem.\\
Romero--Pérez: no infinite descent &
\texttt{well\_founded\_on} in Isabelle/HOL &
The stronger global HOL explication of the source's target-oriented
no-infinite-descent premise; its deletion is witnessed by the integer
predecessor model.\\
\bottomrule
\end{tabular}
\end{center}

\subsection{Cohoe-style essential ordering}

The definition \texttt{essentially\_ordered\_series} requires, for a chosen
target, an explicit derivative predicate and an immediate dependency
predecessor.  More importantly, on the selected ancestry it proves the
equivalence
\[
\texttt{derivative}(x)\ \longleftrightarrow\
\texttt{grounded\_in}(C,D,x).
\]
This equivalence is the development's explicit formal explication of one
reading of derivative efficacy; it is not asserted as a metaphysical identity
between every sense of being caused and every sense of possessing causal
power.  Under that explication, the first member of a founded series is
formally non-derivative rather than merely receiving a free source annotation.
The predicate is still not a synonym for well-foundedness: the integer
predecessor model satisfies the essential-ordering interface while remaining
non-well-founded.  The locale
\texttt{cohoe\_second\_way\_reconstruction} adds the formal translation
conditions needed to obtain the exact conclusion: independent foundation,
singleton endpoint, nonemptiness, and full-ancestry bridges.  These are formal
explications of the source claims, not literal quotations:

\[
\text{a caused target has a non-derivative first cause in the selected carrier}.
\]

The conclusion is intentionally not \lq\lq God exists.''  It is a first-cause
statement in the typed relation-and-carrier interface.

\subsection{An auxiliary bounded stress test for local causal efficacy}

As an auxiliary stress test, the development introduces a deliberately
explicit local support interface.  It requires an immediate cause to exercise
a power possessed by an earlier supporter.  On the integer predecessor model,
the possession and exercise predicates are instantiated by the constantly true
predicates.  The resulting theorem shows that this precise local
support/transmission condition is insufficient to exclude an infinite regress,
even alongside acyclicity and the formal essential-ordering interface.  It is
not a substantive semantics of derivative causal efficacy, nor a claim that
every possible metaphysical theory fails to exclude such a regress; its role is
to mark the exact boundary of this deliberately weak interface.

\subsection{Romero--Pérez reconstruction}

The locale \texttt{romero\_perez\_second\_way\_reconstruction} formalizes the
following package: at least one caused entity exists; caused entities lie in the
selected carrier; every caused entity has a carrier predecessor; and the
carrier-restricted relation is well-founded.  We represent the source's
prohibition of an infinite regress resulting in a particular item by the
stronger global HOL condition \texttt{well\_founded\_on}, rather than by a
target-relative sentence.  The phrase \lq\lq first-order'' describes the source
presentation, not this Isabelle/HOL encoding: \texttt{well\_founded\_on}
quantifies over carrier subsets.  Its exact headline conclusion is

\[
\exists x\in C.\ \texttt{independent\_in}\ C\ D\ x.
\]

For a specified caused target the locale also derives a first cause.  It does
not add uniqueness, necessity, a cause of the universe, or a theological
identity principle.  Its conditions are source-oriented formal assumptions;
the separate premise theory supplies the one-premise-at-a-time countermodel
atlas.

\subsection{Premise factorisation and deletion certificates}

\texttt{Premise\_Minimality.thy} factors the Romero--Pérez formal package into
four conditions: existence, carrier membership, predecessor coverage, and
well-foundedness.  Four deletion theorems retain the others and refute one
fixed combined conclusion: a global independent-support set together with a
first cause for every caused item.  The empty, finite, and integer models
delete these conditions one at a time.

The Cohoe-style translation package uses existence, the formal
essential-ordering explication, full foundation, and three bridges from full
grounding to binary dependence.  Isabelle proves that a separate carrier
condition is redundant because the explication is carrier-restrictive.  Six
deletion theorems refute the formal non-derivative conclusion with small,
inspectable models.  These are minimality results for the encoding, not claims
that the literature has only one interpretation of a source claim.

\section{Related work and limitations}

The formal motivation is Aquinas's \emph{Summa Theologiae}~\cite{AquinasST}.
Cohoe and Romero--Pérez motivate the two source-oriented reconstructions, while
Dixon, Pearson, Bliss, and Rabin--Rabern motivate the distinctions between
grounding, termination, and vicious regress
\cite{Cohoe2013,RomeroPerez2012,Dixon2016,Pearson2023,Bliss2014,RabinRabern2016}.
Prior mechanized philosophical theology provides a precedent for separating a
checked logical core from interpretation~\cite{BenzmuellerPaleo2013}.

The formal contribution is comparative rather than a claim to replace these
source discussions.  The cited philosophical analyses do not supply the
carrier-restricted Isabelle/HOL lattice proved here, its finite and infinite
separation models, or executable certificates whose payloads are connected to
the declarative predicates by kernel-checked theorems.  Cohoe's essentially
ordered series and Romero--Pérez's no-infinite-descent premise are therefore
used as distinct source-oriented inputs to formal explications; the additional
carrier, ancestry, endpoint, and grounding bridges are exposed as translation
assumptions.  The result is a mechanically inspectable comparison of what
follows from those choices, not a claim that the formal vocabulary exhausts
the metaphysics of causation or grounding.

The entry does not formalize empirical causal truth, Aristotelian act and
potency, modal semantics, or the theological identity of a first cause.  It
formalizes consequences of stated relation-theoretic conditions and exhibits
countermodels when a condition is weakened.  All definitions, theorem
statements, proofs, executable examples, and explanatory material needed to
understand the contribution are contained in this AFP entry.

\subsection{Logical background}

The development is carried out in classical Isabelle/HOL.  In particular, the
standard library equivalence between well-foundedness and the absence of an
infinite descending chain is used in
\texttt{well\_founded\_iff\_chain\_terminating}; the encoding is therefore not a
first-order theory and does not claim first-order completeness.  The theorem
statements expose the relevant set and sequence quantification instead of
concealing it behind the phrase \lq\lq no infinite regress''.

\section{Reproducibility}

The Isabelle session is built with
\[
\texttt{isabelle build -D . Thomistic\_Causal\_Arguments}.
\]
For proof-only iteration, use
\texttt{isabelle build -o document=false -D . Thomistic\_Causal\_Arguments}.
The repository-level audit additionally checks the AFP layout, forbidden proof
tokens, stray generated artifacts, and headline theorem counts.  The
Blueprint file records the dependency DAG and the proved foundation/Second-Way
core.

\section{Conclusion}

The result is deliberately narrower than a theological proof: once the
relation, carrier, and foundation conditions are explicit, Isabelle verifies
the hierarchy, certificates, countermodels, and source distinctions.  That
combination is the contribution of the entry.

\tableofcontents

\sloppy
\input{session}
\fussy

\bibliographystyle{abbrv}
\bibliography{root}

\end{document}
